\chapter{Coupled cluster theory}
\label{chap:cctheory}
%=============================================================================

Three defects have followed us through the last five chapters.  Truncated
configuration interaction, Section~\ref{sec:sizeconsistency}, is variational
but not size consistent, so it becomes useless for large systems.
Perturbation theory, Chapter~\ref{chap:mbpt}, is size extensive but not
variational, and its series need not converge.  The random-phase
approximation, Chapter~\ref{chap:mfapplications}, sums one class of
contributions to all orders but collapses when the mean field goes unstable.

Coupled-cluster theory repairs the first two at once.  It is size extensive
by construction, it sums infinite classes of contributions, it is
systematically improvable by raising the excitation rank, and its cost grows
as a modest power of the basis size rather than exponentially.  It is not
variational -- that is the price -- but in practice this matters far less than
one might fear.  It is, by a wide margin, the most successful many-body method
in use: nuclei up to lead have been computed with it, and in quantum chemistry
it is the standard against which everything else is measured.

The idea is a single change of ansatz.  Where configuration interaction writes
$\bket{\Psi}=(1+\hat{C})\bket{\Phi_0}$, coupled cluster writes
$\bket{\Psi}=e^{\hat{T}}\bket{\Phi_0}$.  Everything follows from that
exponential, and this chapter follows it through: the similarity-transformed
Hamiltonian, the reason its expansion terminates exactly, the projected
equations, and then full derivations of the doubles and the singles-and-doubles
amplitude equations.  We work algebraically throughout, as in
Chapter~\ref{chap:mbpt}; the diagrammatic language is a separate subject.  The
pairing model of Section~\ref{sec:pairingmodel} serves as the worked example,
and we compare against every method the book has developed.

%----------------------------------------------------------------------
\section{The exponential ansatz}
\label{sec:exponentialansatz}

\index{coupled cluster!ansatz}\index{cluster operator}
We start, as always, from a reference determinant
\begin{equation}
  \bket{\Phi_0} = \prod_{i\le F}\hat{a}^{\dagger}_i\bket{0},
  \label{eq:10-reference}
\end{equation}
usually but not necessarily the Hartree-Fock state of
Chapter~\ref{chap:hartrefock}.  The exact state is written
\begin{equation}
  \boxed{\;
  \bket{\Psi} = e^{\hat{T}}\bket{\Phi_0},
  \;}
  \label{eq:10-ansatz}
\end{equation}
with the \emph{cluster operator}
\begin{equation}
  \hat{T} = \hat{T}_1 + \hat{T}_2 + \cdots + \hat{T}_N,
  \label{eq:10-clusteroperator}
\end{equation}
\begin{align}
  \hat{T}_1 &= \sum_{ia}t_i^a\,\hat{a}^{\dagger}_a\hat{a}_i,
  \label{eq:10-T1}\\
  \hat{T}_2 &= \frac{1}{4}\sum_{ijab}t_{ij}^{ab}\,
    \hat{a}^{\dagger}_a\hat{a}^{\dagger}_b\hat{a}_j\hat{a}_i,
  \label{eq:10-T2}\\
  \hat{T}_3 &= \frac{1}{36}\sum_{ijkabc}t_{ijk}^{abc}\,
    \hat{a}^{\dagger}_a\hat{a}^{\dagger}_b\hat{a}^{\dagger}_c
    \hat{a}_k\hat{a}_j\hat{a}_i,
  \label{eq:10-T3}
\end{align}
and so on.  The factors $1/4$ and $1/36$ compensate the repeated counting of
antisymmetric index pairs, $(2!)^2$ and $(3!)^2$.  The amplitudes are
antisymmetric under interchange of the occupied or of the virtual indices.

Two features of $\hat{T}$ decide everything that follows.

\paragraph{The cluster operators commute.}
Every term of $\hat{T}$ consists of particle creation and hole annihilation
operators only -- $\hat{a}^{\dagger}_a$ with $a>F$ and $\hat{a}_i$ with
$i\le F$ -- and never the reverse.  The same argument as in
Eq.~(\ref{eq:6-phcommute}) therefore gives
\begin{equation}
  \left[\hat{T}_m,\hat{T}_n\right] = 0
  \label{eq:10-Tcommute}
\end{equation}
for all $m,n$.  This is what makes the exponential tractable.

\paragraph{Excitations cannot be undone.}
Because $\hat{T}$ only excites, $\hat{T}^{n}$ produces at least $n$
particle-hole pairs, and acting on a state with $N$ particles it must vanish
for large enough $n$.  The exponential is therefore a finite sum in any given
matrix element.

\paragraph{Why exponentiate at all?}
Expanding Eq.~(\ref{eq:10-ansatz}),
\begin{equation}
  e^{\hat{T}} = 1 + \hat{T}_1 + \left(\hat{T}_2
    + \tfrac12\hat{T}_1^2\right)
    + \left(\hat{T}_3 + \hat{T}_1\hat{T}_2
      + \tfrac16\hat{T}_1^3\right) + \cdots,
  \label{eq:10-expansion}
\end{equation}
we see that the coefficient of a $2p$-$2h$ determinant is
$t_{ij}^{ab}+t_i^at_j^b-t_i^bt_j^a$, and so on.  Comparing with the
configuration-interaction expansion
$\bket{\Psi}=(1+\hat{C}_1+\hat{C}_2+\cdots)\bket{\Phi_0}$ of
Eq.~(\ref{eq:5-ciexpansion}), the two are related by
\begin{equation}
  \hat{C}_1 = \hat{T}_1,
  \qquad
  \hat{C}_2 = \hat{T}_2 + \tfrac12\hat{T}_1^2,
  \qquad
  \hat{C}_3 = \hat{T}_3 + \hat{T}_1\hat{T}_2 + \tfrac16\hat{T}_1^3,
  \label{eq:10-CtoT}
\end{equation}
and so forth.  The exponential does not add anything the linear expansion
cannot represent; the two are equivalent if both are complete.  What it does
is \emph{reorganise}: the high-rank coefficients are built from products of
low-rank ones rather than treated as independent.  When $\hat{T}$ is
truncated, that reorganisation is exactly what survives.

This is the point.  Truncating $\hat{C}$ at doubles discards the quadruple
excitations entirely, and Section~\ref{sec:sizeconsistency} showed that this
destroys size consistency because a product of two doubles on separate
subsystems \emph{is} a quadruple.  Truncating $\hat{T}$ at doubles keeps that
quadruple, as $\tfrac12\hat{T}_2^2$, with its amplitude fixed by the doubles
rather than free.  Nothing else needs to be said: the disease of
Chapter~\ref{chap:fci} is cured by construction.

%----------------------------------------------------------------------
\section{The similarity-transformed Hamiltonian}
\label{sec:similaritytransform}

\index{similarity transformation}
Insert Eq.~(\ref{eq:10-ansatz}) into the Schr\"odinger equation
$\hat{H}\bket{\Psi}=E\bket{\Psi}$ and multiply from the left by
$e^{-\hat{T}}$:
\begin{equation}
  \boxed{\;
  \overline{H}\bket{\Phi_0} = E\bket{\Phi_0},
  \qquad
  \overline{H} \equiv e^{-\hat{T}}\hat{H}e^{\hat{T}} .
  \;}
  \label{eq:10-transformed}
\end{equation}
This is worth pausing on.  A similarity transformation does not change the
spectrum -- if $\hat{H}\bket{E}=E\bket{E}$ then
$\overline{H}\,e^{-\hat{T}}\bket{E}=E\,e^{-\hat{T}}\bket{E}$ -- so
$\overline{H}$ has exactly the same eigenvalues as $\hat{H}$.  What the
transformation does is move the correlation out of the wave function and into
the operator.  Coupled cluster is the search for a transformation that makes
the \emph{uncorrelated} reference an exact eigenstate.

Because $e^{\hat{T}}$ is not unitary -- $\hat{T}$ is not anti-Hermitian --
$\overline{H}$ is not Hermitian.  That is a genuine cost: left and right
eigenvectors differ, and expectation values require both.  It also brings a
large benefit, as we are about to see.

Normal-ordering with respect to the reference, Section~\ref{sec:normalordering},
splits the Hamiltonian as $\hat{H}=E_{\mathrm{ref}}+\hat{H}_N$ with
\begin{equation}
  \hat{H}_N = \sum_{pq}f_{pq}\left\{\hat{a}^{\dagger}_p\hat{a}_q\right\}
    + \frac{1}{4}\sum_{pqrs}\langle pq\Vert rs\rangle
      \left\{\hat{a}^{\dagger}_p\hat{a}^{\dagger}_q
        \hat{a}_s\hat{a}_r\right\},
  \label{eq:10-HN}
\end{equation}
where $f_{pq}$ is the Fock matrix of Eq.~(\ref{eq:6-fockmatrix}) and
$\langle pq\Vert rs\rangle$ the antisymmetrised two-body element.  Since
$E_{\mathrm{ref}}$ is a number it commutes with $\hat{T}$, so
\begin{equation}
  \overline{H} = E_{\mathrm{ref}} + \overline{H}_N,
  \qquad
  \overline{H}_N = e^{-\hat{T}}\hat{H}_Ne^{\hat{T}},
  \label{eq:10-HbarN}
\end{equation}
and Eq.~(\ref{eq:10-transformed}) becomes
$\overline{H}_N\bket{\Phi_0}=\Delta E\bket{\Phi_0}$ with
$\Delta E=E-E_{\mathrm{ref}}$ the correlation energy.

%----------------------------------------------------------------------
\section{Why the expansion terminates}
\label{sec:termination}

\index{Baker-Campbell-Hausdorff formula!in coupled cluster}
The Baker-Hausdorff lemma of Exercise~\ref{ex:6-bch}d) gives
\begin{equation}
  \overline{H}_N = \hat{H}_N + \left[\hat{H}_N,\hat{T}\right]
    + \frac{1}{2!}\left[\left[\hat{H}_N,\hat{T}\right],\hat{T}\right]
    + \frac{1}{3!}\left[\left[\left[\hat{H}_N,\hat{T}\right],
      \hat{T}\right],\hat{T}\right] + \cdots
  \label{eq:10-bch}
\end{equation}
For a general operator this is an infinite series.  Here it is not, and the
reason is Eq.~(\ref{eq:10-Tcommute}).

Because the cluster operators commute with one another, no commutator in
Eq.~(\ref{eq:10-bch}) can be non-zero by a $\hat{T}$ acting on another
$\hat{T}$: every $\hat{T}$ must contract with $\hat{H}_N$ directly.  In the
language of Section~\ref{sec:contractions}, each $\hat{T}$ must be
\emph{connected} to $\hat{H}_N$.  But $\hat{H}_N$ contains at most four
operators that can be contracted -- two annihilation operators and two
creation operators in its two-body part -- so at most four cluster operators
can attach to it.  Hence
\begin{equation}
  \boxed{\;
  \overline{H}_N = \hat{H}_N + \left[\hat{H}_N,\hat{T}\right]
    + \frac{1}{2!}\left[\left[\hat{H}_N,\hat{T}\right],\hat{T}\right]
    + \frac{1}{3!}\left[\cdots\right]
    + \frac{1}{4!}\left[\cdots\right],
  \;}
  \label{eq:10-bchfinite}
\end{equation}
and the series stops \emph{exactly} at the fourfold commutator.  Every term
is connected, which is why one writes $\overline{H}_N=(\hat{H}_Ne^{\hat{T}})_c$
with the subscript denoting connected terms only.

This is the technical heart of the method, and it is worth appreciating what
has been bought.  The non-Hermiticity of $\overline{H}_N$, which looked like
a defect, is precisely what makes the expansion finite: a unitary
transformation with anti-Hermitian $\hat{T}$ would contain de-excitation
operators, those would not commute among themselves, and the series would
never terminate.  Unitary coupled-cluster theory, which we meet again in the
quantum-computing chapters, pays exactly that price.

%----------------------------------------------------------------------
\section{The coupled-cluster equations}
\label{sec:ccequations}

Equation~(\ref{eq:10-transformed}) says $\bket{\Phi_0}$ is an eigenstate of
$\overline{H}$.  Projecting it onto the reference and onto each excited
determinant turns that statement into a computable set of conditions:
\begin{align}
  \Delta E &= \bracket{\Phi_0}{\overline{H}_N}{\Phi_0},
  \label{eq:10-energyeq}\\
  0 &= \bracket{\Phi_i^a}{\overline{H}_N}{\Phi_0},
  \label{eq:10-singleseq}\\
  0 &= \bracket{\Phi_{ij}^{ab}}{\overline{H}_N}{\Phi_0},
  \label{eq:10-doubleseq}
\end{align}
and so on up the excitation ladder.  The first line gives the energy once the
amplitudes are known; the rest determine them.

Read physically, Eqs.~(\ref{eq:10-singleseq}) and~(\ref{eq:10-doubleseq}) say
that the similarity-transformed Hamiltonian has \emph{no} matrix elements
connecting the reference to $1p$-$1h$ or $2p$-$2h$ states.  The reference has
been decoupled from the simple excitations.  It has not been decoupled from
everything: with $\hat{T}=\hat{T}_1+\hat{T}_2$, the operator $\overline{H}_N$
still connects $\bket{\Phi_0}$ to $3p$-$3h$ and $4p$-$4h$ states, and that
residual coupling is the entire error of CCSD.

\begin{notebox}
\textbf{Why is CCD not exact?}  A similarity transformation preserves all
eigenvalues, so one might expect the transformed problem to be exactly
soluble.  The resolution is that we are not solving the transformed problem
exactly either.  With $\hat{T}=\hat{T}_2$ the full $\overline{H}_N$ contains
up to six-body terms -- $\hat{H}_N$ has four contractible operators, and each
$\hat{T}_2$ adds two more -- and it is that full operator which would
reproduce the exact energy.  What we impose is only that the $2p$-$2h$ block
of $\overline{H}_N$ vanishes.  Coupled cluster at any finite truncation is
therefore a similarity transformation \emph{plus} a projection, and the
projection is where the approximation lives.
\end{notebox}

%----------------------------------------------------------------------
\section{The correlation energy}
\label{sec:ccenergy}

Evaluating Eq.~(\ref{eq:10-energyeq}) is the easiest part of the calculation.
Only terms of $\overline{H}_N$ that de-excite by exactly the amount
$\hat{T}$ excites can contribute, so at most two cluster operators can
attach, and
\begin{equation}
  \Delta E = \bracket{\Phi_0}{\hat{H}_N\left(1+\hat{T}_1+\hat{T}_2
    +\tfrac12\hat{T}_1^2\right)}{\Phi_0}_c .
  \label{eq:10-energyterms}
\end{equation}
The one-body part of $\hat{H}_N$ contracts with $\hat{T}_1$, the two-body part
with $\hat{T}_2$ and with $\tfrac12\hat{T}_1^2$, and nothing else survives:
\begin{equation}
  \boxed{\;
  \Delta E = \sum_{ia}f_{ia}\,t_i^a
    + \frac{1}{4}\sum_{ijab}\langle ij\Vert ab\rangle\,t_{ij}^{ab}
    + \frac{1}{2}\sum_{ijab}\langle ij\Vert ab\rangle\,t_i^at_j^b .
  \;}
  \label{eq:10-ccsdenergy}
\end{equation}
In a Hartree-Fock basis $f_{ia}=0$ by Brillouin's theorem,
Eq.~(\ref{eq:6-brillouin}), and the first term disappears.  The singles then
enter the energy only through the last term, quadratically -- which is why
CCSD energies are often close to CCD energies even when the singles
amplitudes are not small.

Notice also that Eq.~(\ref{eq:10-ccsdenergy}) has exactly the form of the
configuration-interaction correlation energy, Eq.~(\ref{eq:9-fciDE}), with
$C_{ij}^{ab}$ replaced by $t_{ij}^{ab}+t_i^at_j^b-t_i^bt_j^a$.  That
combination will appear so often that it is given a name,
\begin{equation}
  \tau_{ij}^{ab} = t_{ij}^{ab} + t_i^at_j^b - t_i^bt_j^a,
  \label{eq:10-tau}
\end{equation}
and its half-weighted cousin
\begin{equation}
  \tilde{\tau}_{ij}^{ab} = t_{ij}^{ab}
    + \tfrac12\left(t_i^at_j^b - t_i^bt_j^a\right).
  \label{eq:10-tautilde}
\end{equation}

%----------------------------------------------------------------------
\section{Deriving the CCD equations}
\label{sec:ccdderivation}

\index{coupled cluster!doubles}
We derive the doubles equation first, with $\hat{T}=\hat{T}_2$ alone.  This is
the \emph{coupled-cluster doubles} approximation, and it is exact for any
system whose interaction cannot produce a $1p$-$1h$ excitation -- the pairing
model, nuclear matter on a momentum grid -- and the leading approximation in
a Hartree-Fock basis for everything else.

The equation to satisfy is
$\bracket{\Phi_{ij}^{ab}}{\overline{H}_N}{\Phi_0}=0$ with
\begin{equation}
  \overline{H}_N = \left(\hat{H}_N e^{\hat{T}_2}\right)_c
   = \left(\hat{H}_N\left[1+\hat{T}_2
     + \tfrac12\hat{T}_2^2\right]\right)_c .
  \label{eq:10-ccdexpansion}
\end{equation}
The series stops at $\hat{T}_2^2$: a third $\hat{T}_2$ would need two more
contractible operators in $\hat{H}_N$ than it has, once the two free particle
and two free hole lines of the $2p$-$2h$ projection are accounted for.

We take the three terms in turn.  Throughout we use the permutation
operators
\begin{equation}
  P(ij)X_{ij} = X_{ij}-X_{ji},
  \qquad
  P(ab)X^{ab} = X^{ab}-X^{ba},
  \label{eq:10-permutation}
\end{equation}
and their product
$P(ij)P(ab)X_{ij}^{ab}=X_{ij}^{ab}-X_{ji}^{ab}-X_{ij}^{ba}+X_{ji}^{ba}$,
which enforce the antisymmetry of the amplitudes under interchange of the
two occupied or the two virtual labels.  Repeated indices are summed.

\paragraph{No cluster operator.}
$\bracket{\Phi_{ij}^{ab}}{\hat{H}_N}{\Phi_0}$ is the Condon-Slater matrix
element between the reference and a double excitation,
Section~\ref{sec:wickapps}, which is simply
\begin{equation}
  \langle ab\Vert ij\rangle .
  \label{eq:10-ccd0}
\end{equation}
This is the driving term: without it all amplitudes would vanish.

\paragraph{One cluster operator.}
$\bracket{\Phi_{ij}^{ab}}{\left(\hat{H}_N\hat{T}_2\right)_c}{\Phi_0}$
requires all four external lines of the projection to be accounted for, with
$\hat{T}_2$ supplying two particle and two hole lines and $\hat{H}_N$
connecting to the rest.  Enumerating the ways to contract, the one-body part
of $\hat{H}_N$ can attach to a particle line or to a hole line,
\begin{equation}
  P(ab)\sum_c f_{bc}\,t_{ij}^{ac}
  - P(ij)\sum_k f_{kj}\,t_{ik}^{ab},
  \label{eq:10-ccd1a}
\end{equation}
and the two-body part in three distinct ways -- both particle lines, both hole
lines, or one of each:
\begin{equation}
  \frac{1}{2}\sum_{cd}\langle ab\Vert cd\rangle\,t_{ij}^{cd}
  + \frac{1}{2}\sum_{kl}\langle kl\Vert ij\rangle\,t_{kl}^{ab}
  + P(ij)P(ab)\sum_{kc}\langle kb\Vert cj\rangle\,t_{ik}^{ac}.
  \label{eq:10-ccd1b}
\end{equation}
These three are the \emph{particle-particle ladder}, the \emph{hole-hole
ladder} and the \emph{ring} term.  The factors $\tfrac12$ compensate the
double counting of the two equivalent contractions in each ladder; the ring
has no such factor because the two lines are inequivalent, and instead
carries the full permutation $P(ij)P(ab)$.

\paragraph{Two cluster operators.}
$\tfrac12\bracket{\Phi_{ij}^{ab}}{\left(\hat{H}_N\hat{T}_2^2\right)_c}
{\Phi_0}$ needs both $\hat{T}_2$ connected to $\hat{H}_N$, which forces the
two-body part and uses all four of its contractible operators.  Only
$\langle kl\Vert cd\rangle$ with two hole and two particle indices can do
this.  There are four distinct topologies, according to how the four
contractions are distributed between the two amplitudes:
\begin{align}
  &+\frac{1}{4}\sum_{klcd}\langle kl\Vert cd\rangle\,
     t_{ij}^{cd}\,t_{kl}^{ab}
  \nonumber\\
  &+ P(ij)P(ab)\,\frac{1}{2}\sum_{klcd}\langle kl\Vert cd\rangle\,
     t_{ik}^{ac}\,t_{jl}^{bd}
  \nonumber\\
  &- P(ij)\,\frac{1}{2}\sum_{klcd}\langle kl\Vert cd\rangle\,
     t_{ik}^{cd}\,t_{lj}^{ab}
  \nonumber\\
  &- P(ab)\,\frac{1}{2}\sum_{klcd}\langle kl\Vert cd\rangle\,
     t_{kl}^{ac}\,t_{ij}^{db}.
  \label{eq:10-ccd2}
\end{align}
The first is the product of a hole-hole and a particle-particle ladder; the
second is a ring on each side; the third and fourth attach both indices of
one amplitude to the interaction and leave the other amplitude as a spectator.

\paragraph{The result.}
Collecting Eqs.~(\ref{eq:10-ccd0})--(\ref{eq:10-ccd2}) and writing the
diagonal Fock elements on the left, the CCD amplitude equation is
\begin{align}
  D_{ij}^{ab}\,t_{ij}^{ab}
  &= \langle ab\Vert ij\rangle
   + P(ab)\sum_c f_{bc}^{\,\prime}\,t_{ij}^{ac}
   - P(ij)\sum_k f_{kj}^{\,\prime}\,t_{ik}^{ab}
  \nonumber\\
  &\quad
   + \frac{1}{2}\sum_{cd}\langle ab\Vert cd\rangle\,t_{ij}^{cd}
   + \frac{1}{2}\sum_{kl}\langle kl\Vert ij\rangle\,t_{kl}^{ab}
   + P(ij)P(ab)\sum_{kc}\langle kb\Vert cj\rangle\,t_{ik}^{ac}
  \nonumber\\
  &\quad
   + \frac{1}{4}\sum_{klcd}\langle kl\Vert cd\rangle\,t_{ij}^{cd}t_{kl}^{ab}
   + \frac{1}{2}P(ij)P(ab)\sum_{klcd}\langle kl\Vert cd\rangle\,
     t_{ik}^{ac}t_{jl}^{bd}
  \nonumber\\
  &\quad
   - \frac{1}{2}P(ij)\sum_{klcd}\langle kl\Vert cd\rangle\,
     t_{ik}^{cd}t_{lj}^{ab}
   - \frac{1}{2}P(ab)\sum_{klcd}\langle kl\Vert cd\rangle\,
     t_{kl}^{ac}t_{ij}^{db},
  \label{eq:10-ccd}
\end{align}
with the energy denominator
\begin{equation}
  D_{ij}^{ab} = f_{ii}+f_{jj}-f_{aa}-f_{bb}
   = \varepsilon_i+\varepsilon_j-\varepsilon_a-\varepsilon_b
  \label{eq:10-denominator}
\end{equation}
and $f'$ the off-diagonal part of the Fock matrix, which vanishes in a
canonical Hartree-Fock basis.

Equation~(\ref{eq:10-ccd}) is a set of coupled quadratic equations, one for
each amplitude, and it is solved by iteration: start from
$t_{ij}^{ab}=\langle ab\Vert ij\rangle/D_{ij}^{ab}$, which is the second-order
perturbation amplitude of Eq.~(\ref{eq:9-firstordercoefficients}), evaluate
the right-hand side, divide by $D$, and repeat.

\paragraph{Intermediates.}
Evaluated naively, the quadratic terms cost $\bigO(n^8)$.  Factoring them
through intermediates that are built once per iteration reduces this to
$\bigO(n^6)$, which is what makes the method usable.  Defining
\begin{align}
  \chi_{kj} &= \frac{1}{2}\sum_{lcd}
    \langle kl\Vert cd\rangle\,t_{jl}^{cd},
  &
  \chi_{cb} &= -\frac{1}{2}\sum_{kld}
    \langle kl\Vert cd\rangle\,t_{kl}^{bd},
  \nonumber\\
  \chi_{klij} &= \frac{1}{2}\sum_{cd}
    \langle kl\Vert cd\rangle\,t_{ij}^{cd},
  &
  \chi_{bkcj} &= \frac{1}{2}\sum_{ld}
    \langle kl\Vert cd\rangle\,t_{lj}^{db},
  \label{eq:10-intermediates}
\end{align}
the last four lines of Eq.~(\ref{eq:10-ccd}) become
\begin{equation}
  -P(ij)\sum_k t_{ik}^{ab}\chi_{kj}
  + P(ab)\sum_c t_{ij}^{ac}\chi_{cb}
  + \frac{1}{2}\sum_{kl}t_{kl}^{ab}\chi_{klij}
  + P(ij)P(ab)\sum_{kc}t_{ik}^{ac}\chi_{bkcj},
  \label{eq:10-intermediateform}
\end{equation}
each a single matrix multiplication.  This is how the program implements it,
and how every production code does.

%----------------------------------------------------------------------
\section{Deriving the CCSD equations}
\label{sec:ccsdderivation}

\index{coupled cluster!singles and doubles}
With $\hat{T}=\hat{T}_1+\hat{T}_2$ the derivation is the same in principle and
far longer in practice, because $\hat{T}_1$ can appear up to four times.  We
give the result in the standard factored form, which is both the clearest
statement and the one that is implemented.

The organising trick is to collect the recurring combinations into
intermediates.  Three carry two indices,
\begin{align}
  F_{ae} &= f_{ae}^{\,\prime}
    - \frac{1}{2}\sum_m f_{me}t_m^a
    + \sum_{mf}t_m^f\langle ma\Vert fe\rangle
    - \frac{1}{2}\sum_{mnf}\tilde{\tau}_{mn}^{af}
      \langle mn\Vert ef\rangle,
  \label{eq:10-Fae}\\
  F_{mi} &= f_{mi}^{\,\prime}
    + \frac{1}{2}\sum_e t_i^e f_{me}
    + \sum_{en}t_n^e\langle mn\Vert ie\rangle
    + \frac{1}{2}\sum_{nef}\tilde{\tau}_{in}^{ef}
      \langle mn\Vert ef\rangle,
  \label{eq:10-Fmi}\\
  F_{me} &= f_{me}
    + \sum_{nf}t_n^f\langle mn\Vert ef\rangle,
  \label{eq:10-Fme}
\end{align}
and three carry four,
\begin{align}
  W_{mnij} &= \langle mn\Vert ij\rangle
    + P(ij)\sum_e t_j^e\langle mn\Vert ie\rangle
    + \frac{1}{4}\sum_{ef}\tau_{ij}^{ef}\langle mn\Vert ef\rangle,
  \label{eq:10-Wmnij}\\
  W_{abef} &= \langle ab\Vert ef\rangle
    - P(ab)\sum_m t_m^b\langle am\Vert ef\rangle
    + \frac{1}{4}\sum_{mn}\tau_{mn}^{ab}\langle mn\Vert ef\rangle,
  \label{eq:10-Wabef}\\
  W_{mbej} &= \langle mb\Vert ej\rangle
    + \sum_f t_j^f\langle mb\Vert ef\rangle
    - \sum_n t_n^b\langle mn\Vert ej\rangle
    \nonumber\\
  &\quad
    - \sum_{nf}\left(\tfrac12 t_{jn}^{fb} + t_j^ft_n^b\right)
      \langle mn\Vert ef\rangle .
  \label{eq:10-Wmbej}
\end{align}
These are the intermediates of Stanton, Gauss, Watts and Bartlett, and their
structure repays a moment's attention: each is a bare integral, plus
corrections in which the amplitudes dress it.  $F_{ae}$ and $F_{mi}$ are the
Fock matrix dressed by correlation; $W_{mnij}$ and $W_{abef}$ are the
hole-hole and particle-particle ladders dressed likewise; and $W_{mbej}$ is
the dressed ring.

With them, the amplitude equations take a compact form.  The singles:
\begin{align}
  D_i^a\,t_i^a
  &= f_{ia}
   + \sum_e t_i^e F_{ae}
   - \sum_m t_m^a F_{mi}
   + \sum_{me}t_{im}^{ae}F_{me}
  \nonumber\\
  &\quad
   - \sum_{nf}t_n^f\langle na\Vert if\rangle
   - \frac{1}{2}\sum_{mef}t_{im}^{ef}\langle ma\Vert ef\rangle
   - \frac{1}{2}\sum_{men}t_{mn}^{ae}\langle nm\Vert ei\rangle,
  \label{eq:10-ccsdsingles}
\end{align}
with $D_i^a=\varepsilon_i-\varepsilon_a$.  The doubles:
\begin{align}
  D_{ij}^{ab}\,t_{ij}^{ab}
  &= \langle ij\Vert ab\rangle
   + P(ab)\sum_e t_{ij}^{ae}
     \left(F_{be}-\tfrac12\sum_m t_m^b F_{me}\right)
  \nonumber\\
  &\quad
   - P(ij)\sum_m t_{im}^{ab}
     \left(F_{mj}+\tfrac12\sum_e t_j^e F_{me}\right)
  \nonumber\\
  &\quad
   + \frac{1}{2}\sum_{mn}\tau_{mn}^{ab}W_{mnij}
   + \frac{1}{2}\sum_{ef}\tau_{ij}^{ef}W_{abef}
  \nonumber\\
  &\quad
   + P(ij)P(ab)\sum_{me}\left(t_{im}^{ae}W_{mbej}
     - t_i^et_m^a\langle mb\Vert ej\rangle\right)
  \nonumber\\
  &\quad
   + P(ij)\sum_e t_i^e\langle ab\Vert ej\rangle
   - P(ab)\sum_m t_m^a\langle mb\Vert ij\rangle .
  \label{eq:10-ccsddoubles}
\end{align}

\paragraph{Reading the equations.}
Several structural features are worth noting, and each connects to an earlier
chapter.

\emph{Setting all amplitudes on the right to zero} leaves
$D_{ij}^{ab}t_{ij}^{ab}=\langle ij\Vert ab\rangle$, which is the first-order
perturbation amplitude, and $D_i^at_i^a=f_{ia}$, which vanishes in a
Hartree-Fock basis.  The first iteration of a coupled-cluster calculation is
therefore an MP2 calculation, exactly as
Eq.~(\ref{eq:9-firstordercoefficients}) said.

\emph{Setting $\hat{T}_1=0$} makes $\tau=\tilde\tau=t_2$, kills the last two
lines of Eq.~(\ref{eq:10-ccsddoubles}) and reduces the intermediates to those
of Eq.~(\ref{eq:10-intermediates}); Eq.~(\ref{eq:10-ccd}) is recovered.

\emph{The singles equation is not driven} in a Hartree-Fock basis: with
$f_{ia}=0$ the leading surviving term is
$-\tfrac12\sum_{mef}t_{im}^{ef}\langle ma\Vert ef\rangle$, which is first
order in $t_2$ and hence second order in the interaction.
Singles therefore start one order later than doubles, which is Brillouin's
theorem showing up again.

\emph{The equations are nonlinear} in both $t_1$ and $t_2$, and it is the
nonlinearity that performs the infinite resummations.  The
Section~\ref{sec:ccdcontent} makes that statement quantitative.

\begin{notebox}
\textbf{Cost.}  The most expensive term in Eq.~(\ref{eq:10-ccsddoubles}) is
the particle-particle ladder, $\tfrac12\tau_{ij}^{ef}W_{abef}$, which costs
$\bigO(n_o^2n_v^4)$.  With $n_o$ occupied and $n_v$ virtual orbitals the
total scales as $n^6$, against the exponential scaling of full configuration
interaction, Section~\ref{sec:exponentialwall}.  This single fact is why
coupled cluster reaches systems that configuration interaction cannot
approach: for a hundred orbitals, $n^6=10^{12}$ operations is a large but
finite calculation, while $\binom{100}{50}\approx10^{29}$ determinants is not.
Adding triples, CCSDT, costs $n^8$; the perturbative CCSD(T) approximation
costs $n^7$ and is the standard compromise.
\end{notebox}

%----------------------------------------------------------------------
\section{Coupled cluster as a resummation of perturbation theory}
\label{sec:ccdcontent}

\index{coupled cluster!and perturbation theory}
Chapter~\ref{chap:mbpt} built the correlation energy as a series in the
interaction and found that it may converge slowly or not at all.  Coupled
cluster contains the same series, but does not sum it term by term.  The
relation can be made exact and, for the pairing model, checked numerically to
every digit.

The CCD amplitude equation~(\ref{eq:10-ccd}) is \emph{quadratic} in the
amplitudes -- nothing of higher degree appears -- so it may be written
\begin{equation}
  D\,t = C + L[t] + Q[t,t],
  \label{eq:10-quadratic}
\end{equation}
with $C$ the driving integral, $L$ the linear terms and $Q$ the bilinear ones.
Expanding $t=\sum_k t^{(k)}$ with $t^{(k)}$ of order $k$ in the interaction and
collecting powers,
\begin{equation}
  D\,t^{(1)} = C,
  \qquad
  D\,t^{(k)} = L\!\left[t^{(k-1)}\right]
    + \sum_{i+j=k-1}Q\!\left[t^{(i)},t^{(j)}\right],
  \label{eq:10-orderbyorder}
\end{equation}
since $L$ and $Q$ each carry one power of the interaction.  The energy
contribution of $t^{(k)}$ is of order $k+1$.

This is a perturbation expansion, and it is the M\o{}ller-Plesset one: the
first amplitude is $\langle ab\Vert ij\rangle/D_{ij}^{ab}$ and the first
energy is MP2.  The program checks the identification against the
Rayleigh-Schr\"odinger machinery of Chapter~\ref{chap:mbpt}, computed
independently in the determinant basis:

\begin{center}
\renewcommand{\arraystretch}{1.2}
\setlength{\tabcolsep}{11pt}
\begin{tabular}{rrrr}
\toprule
$g$ & order & from CCD & from MBPT\\
\midrule
$0.50$ & $2$ & $-0.0623931624$ & $-0.0623931624$\\
       & $3$ & $-0.0165183724$ & $-0.0165183724$\\
       & $4$ & $-0.0038210766$ & $-0.0038210766$\\
$1.00$ & $2$ & $-0.2190476190$ & $-0.2190476190$\\
       & $3$ & $-0.1004988662$ & $-0.1004988662$\\
       & $4$ & $-0.0403096858$ & $-0.0403096858$\\
\bottomrule
\end{tabular}
\end{center}

They agree to every digit through fourth order.  That CCD reproduces MP2 and
MP3 exactly is general -- those orders contain only doubles.  That it also
gets MP4 right here is special to the pairing model, where the interaction
produces neither singles nor triples and the quadruple excitations enter only
through the disconnected $\tfrac12\hat{T}_2^2$ that the exponential already
supplies.  For a general interaction, MP4 contains connected triples that CCD
does not have, and the agreement stops at third order.

\paragraph{A warning about partitions.}
The numbers above are \emph{not} those of Chapter~\ref{chap:mbpt}, and the
difference is instructive.  There we took $\hat{H}_0$ to be the bare one-body
operator; here the denominators contain the Fock matrix, whose occupied
elements are shifted down by $g/2$.  For the pairing model the second-order
energy in the two partitions is
\begin{equation}
  \Delta E^{(2)}_{\text{one-body}} = -\frac{7}{24}g^2,
  \qquad
  \Delta E^{(2)}_{\text{M\o{}ller-Plesset}}
   = -\frac{g^2}{4}\left[\frac{1}{2+g}+\frac{2}{4+g}
     +\frac{1}{6+g}\right],
  \label{eq:10-twopartitions}
\end{equation}
which differ already at order $g^3$: at $g=1$ they give $-0.2917$ and
$-0.2190$.  Section~\ref{sec:partitions} warned that a perturbative number
without its partition is meaningless, and here is the warning in force.
Coupled cluster uses the M\o{}ller-Plesset partition because the Fock operator
is what makes the singles vanish at first order.

\paragraph{Summing the series.}
The point of the exponential is that the expansion need not be summed term by
term.  For the pairing model:

\begin{table}[h]
\centering
\renewcommand{\arraystretch}{1.25}
\setlength{\tabcolsep}{9pt}
\small
\begin{tabular}{rrrr}
\toprule
order & $g=0.5$ & $g=1.0$ & $g=1.5$\\
\midrule
$2$ & $-6.24\times10^{-2}$ & $-2.19\times10^{-1}$ & $-4.40\times10^{-1}$\\
$3$ & $-1.65\times10^{-2}$ & $-1.00\times10^{-1}$ & $-2.68\times10^{-1}$\\
$4$ & $-3.82\times10^{-3}$ & $-4.03\times10^{-2}$ & $-1.43\times10^{-1}$\\
$5$ & $-6.58\times10^{-4}$ & $-1.23\times10^{-2}$ & $-5.86\times10^{-2}$\\
$6$ & $-2.98\times10^{-5}$ & $-1.31\times10^{-3}$ & $-9.80\times10^{-3}$\\
$7$ & $+3.49\times10^{-5}$ & $+1.62\times10^{-3}$ & $+1.16\times10^{-2}$\\
\midrule
sum of $16$ & $-0.0833623$ & $-0.3695742$ & $-0.8857526$\\
CCD & $-0.0833623$ & $-0.3695572$ & $-0.8841653$\\
\bottomrule
\end{tabular}
\caption{The order-by-order content of coupled-cluster doubles for the
pairing model, and its resummation.  The expansion converges back onto the
CCD energy; solving the nonlinear equation reaches the same answer directly
and without the intermediate oscillation.}
\label{tab:ccdorders}
\end{table}

The series does converge here -- but slowly, and non-monotonically, changing
sign at seventh order.  Solving Eq.~(\ref{eq:10-ccd}) by iteration reaches
the same number in twenty-odd steps, with no reference to the order
structure at all.  That is the practical content of ``infinite resummation'':
not that new physics appears, but that the physics already present in the
series is obtained without having to sum it.

The classes summed are those the truncation can represent.  The
particle-particle ladder term of Eq.~(\ref{eq:10-ccd}), iterated, sums the
ladder series to all orders; the ring term sums the ring series, which is
what the random-phase approximation of Chapter~\ref{chap:mfapplications} does
in its own channel; and the quadratic terms couple them.  CCD is neither a
ladder theory nor a ring theory but a self-consistent combination of both.

%----------------------------------------------------------------------
\section{The pairing model}
\label{sec:ccpairing}

We now do the calculation.  The pairing Hamiltonian of
Section~\ref{sec:pairingmodel},
\begin{equation}
  \hat{H} = \xi\sum_{p\sigma}(p-1)\,
    \hat{a}^{\dagger}_{p\sigma}\hat{a}_{p\sigma}
  - \frac{g}{2}\sum_{pq}\hat{a}^{\dagger}_{p+}\hat{a}^{\dagger}_{p-}
    \hat{a}_{q-}\hat{a}_{q+},
  \label{eq:10-pairing}
\end{equation}
has $\langle p{+}\,p{-}\Vert q{+}\,q{-}\rangle=-g/2$ and nothing else.  Two
consequences follow immediately, and both simplify the calculation
enormously.

First, the interaction moves complete pairs, so it has no matrix element
between the reference and a $1p$-$1h$ state.  Hence $t_i^a=0$ identically and
CCSD reduces to CCD.  The program verifies this: the singles amplitudes come
out as exactly zero, and the two energies agree to every digit.

Second, only three of the interaction channels are non-zero --
$\langle ab\Vert cd\rangle$, $\langle ab\Vert ij\rangle$ and
$\langle ij\Vert kl\rangle$.  There is no $\langle ab\Vert id\rangle$ or
$\langle ai\Vert jb\rangle$, because those would require breaking a pair.  The
ring term of Eq.~(\ref{eq:10-ccd}) therefore vanishes, and only the two
ladders and the quadratic terms survive.

The Fock matrix is diagonal, with
\begin{equation}
  \varepsilon_i = \xi(p_i-1) - \frac{g}{2},
  \qquad
  \varepsilon_a = \xi(p_a-1),
  \label{eq:10-pairingfock}
\end{equation}
the occupied levels shifted down by $g/2$ and the empty ones untouched --
which is the Hartree-Fock solution found in Section~\ref{sec:hfnothing}.
Taking four levels and four particles, the results are

\begin{table}[h]
\centering
\renewcommand{\arraystretch}{1.25}
\setlength{\tabcolsep}{7pt}
\small
\begin{tabular}{rrrrrrr}
\toprule
$g$ & $E_{\mathrm{ref}}$ & MP2 & CCD & FCI & CCD error & iterations\\
\midrule
$0.25$ & $1.75000$ & $1.73320261$ & $1.73045621$ & $1.73045985$
  & $-3.6\times10^{-6}$ & $13$\\
$0.50$ & $1.50000$ & $1.43760684$ & $1.41663766$ & $1.41677428$
  & $-1.4\times10^{-4}$ & $18$\\
$1.00$ & $1.00000$ & $0.78095238$ & $0.63044275$ & $0.63554847$
  & $-5.1\times10^{-3}$ & $26$\\
$1.50$ & $0.50000$ & $0.05974026$ & $-0.38416526$ & $-0.34819051$
  & $-3.6\times10^{-2}$ & $30$\\
$2.00$ & $0.00000$ & $-0.70833333$ & $-1.60959440$ & $-1.48965216$
  & $-1.2\times10^{-1}$ & $31$\\
\bottomrule
\end{tabular}
\caption{Coupled-cluster doubles for the pairing model with four doubly
degenerate levels and four particles, $\xi=1$.  The reference energy is the
Hartree-Fock energy $2-g$ of Chapter~\ref{chap:hartrefock}, and the exact
column is the full configuration interaction of Chapter~\ref{chap:fci}.}
\label{tab:ccdpairing}
\end{table}

At weak coupling CCD is very nearly exact -- four parts in a million at
$g=0.25$, where MP2 is already wrong in the third decimal.  It
over-correlates, and the error grows with the coupling, reaching $0.12$ at
$g=2$.  What is missing is the \emph{connected} quadruple excitations: the
exponential supplies the $4p$-$4h$ amplitude as $\tfrac12\hat{T}_2^2$, the
disconnected part only, and for four particles in four levels the connected
part is not negligible once the interaction is strong.  Including $\hat{T}_4$
would make the calculation exact, since the exact state has no components
beyond $4p$-$4h$.

%----------------------------------------------------------------------
\section{When singles matter}
\label{sec:ccsinglesmatter}

To see CCSD do something CCD cannot, we need an interaction that breaks
pairs.  The pairing plus particle-hole model of
Section~\ref{sec:pairingphmodel} is exactly that: its particle-hole term
\begin{equation}
  \hat{V}_{\mathrm{ph}} = -\frac{f}{2}\sum_{pqr}
    \left(\hat{a}^{\dagger}_{p+}\hat{a}^{\dagger}_{p-}
      \hat{a}_{q-}\hat{a}_{r+} + \mathrm{h.c.}\right)
  \label{eq:10-phterm}
\end{equation}
connects the reference to $1p$-$1h$ states and the singles amplitudes are no
longer zero.  We run the calculation twice, once on the bare oscillator
reference and once on the Hartree-Fock reference of
Chapter~\ref{chap:hartrefock}, where Brillouin's theorem holds.

\begin{table}[h]
\centering
\renewcommand{\arraystretch}{1.25}
\setlength{\tabcolsep}{5pt}
\small
\begin{tabular}{rrrrrrrr}
\toprule
$g$ & $f$ & CCD & CCSD & CCSD (HF) & FCI
  & $\max|t_1|$ & $\max|t_1|$ (HF)\\
\midrule
$0.50$ & $0.025$ & $1.3476038$ & $1.3469075$ & $1.3469088$ & $1.3471327$
  & $1.1\times10^{-2}$ & $5.4\times10^{-4}$\\
$0.70$ & $0.050$ & $0.9707634$ & $0.9681476$ & $0.9681644$ & $0.9697588$
  & $1.9\times10^{-2}$ & $1.5\times10^{-3}$\\
$1.00$ & $0.050$ & $0.4447429$ & $0.4423659$ & $0.4424119$ & $0.4505823$
  & $1.7\times10^{-2}$ & $1.8\times10^{-3}$\\
$0.50$ & $0.250$ & $0.6033126$ & $0.5467210$ & $0.5472745$ & $0.5511916$
  & $8.4\times10^{-2}$ & $7.3\times10^{-3}$\\
$1.00$ & $0.500$ & $-1.6764657$ & $-1.7957374$ & $-1.7805416$ & $-1.6917367$
  & $1.0\times10^{-1}$ & $1.3\times10^{-2}$\\
\bottomrule
\end{tabular}
\caption{CCD and CCSD for the pairing plus particle-hole model of
Chapter~\ref{chap:mfapplications}, on the bare oscillator reference and on
the Hartree-Fock reference.}
\label{tab:ccsdph}
\end{table}

Three things to read off.

\emph{The singles matter.}  At $g=0.5$, $f=0.25$ they move the energy from
$0.6033$ to $0.5467$ against an exact $0.5512$, turning a nine per cent error
into one of under one per cent.

\emph{Brillouin's theorem is visible.}  The singles amplitudes are an order of
magnitude smaller on the Hartree-Fock reference, because there they start at
second order in the interaction rather than first,
Eq.~(\ref{eq:10-ccsdsingles}).  This is the practical reason for working in a
Hartree-Fock basis: not that the singles vanish, but that they are small.

\emph{Coupled cluster is not invariant under a change of reference.}  The two
CCSD columns are close but not identical.  A truncated cluster operator
depends on the determinant it is built from, and only the untruncated theory
is reference-independent.  In a Hartree-Fock basis, where $\hat{T}_1$ has
little to do, the dependence is weak -- which is the usual situation and the
usual justification.

Note also the last row.  At $g=1$, $f=0.5$ the particle-hole term is strong,
CCSD overshoots the exact energy by $0.11$, and CCD -- which cannot see the
singles at all -- happens to be closer.  This is not a reason to prefer CCD.
It is a reminder that an uncontrolled error can have either sign, and that
agreement at a single point proves nothing.

%----------------------------------------------------------------------
\section{Size extensivity}
\label{sec:ccextensivity}

\index{size extensivity!coupled cluster}
The property that motivated the whole construction can now be verified.  For
two non-interacting subsystems $A$ and $B$, the cluster operator separates,
$\hat{T}=\hat{T}_A+\hat{T}_B$, the two pieces commute, and therefore
\begin{equation}
  e^{\hat{T}}\bket{\Phi_0}
   = e^{\hat{T}_A}e^{\hat{T}_B}\bket{\Phi_A}\bket{\Phi_B}
   = \left(e^{\hat{T}_A}\bket{\Phi_A}\right)
     \left(e^{\hat{T}_B}\bket{\Phi_B}\right).
  \label{eq:10-factorises}
\end{equation}
The wave function factorises exactly, and the energy is therefore additive --
at any truncation level, since $\hat{T}_A$ and $\hat{T}_B$ are truncated
independently.  This is the whole argument; it is one line, and it is the
reason coupled cluster exists.

The program confirms it for identical copies of a three-level pairing system:

\begin{center}
\renewcommand{\arraystretch}{1.2}
\setlength{\tabcolsep}{11pt}
\begin{tabular}{rrrr}
\toprule
$g$ & $E(1)$ & $E(2)-2E(1)$ & $E(3)-3E(1)$\\
\midrule
$0.25$ & $-0.012167069342$ & $0$ & $0$\\
$0.50$ & $-0.050059379847$ & $-2.1\times10^{-13}$ & $-3.2\times10^{-13}$\\
$1.00$ & $-0.205303400091$ & $\phantom{-}5.6\times10^{-17}$
  & $-3.0\times10^{-13}$\\
\bottomrule
\end{tabular}
\end{center}

Exact to machine precision.  Compare Section~\ref{sec:sizeconsistency},
where doubles configuration interaction -- the \emph{linear} theory with the
same excitation rank -- had errors growing from $6\times10^{-5}$ to
$1.1\times10^{-2}$ over the same range of couplings.  The
two methods keep the same excitations; only the ansatz differs; and that
difference decides whether the method can be used on a solid.

%----------------------------------------------------------------------
\section{Everything against everything}
\label{sec:cccomparison}

The pairing model has now been treated by every method in this book.

\begin{table}[h]
\centering
\renewcommand{\arraystretch}{1.25}
\setlength{\tabcolsep}{6pt}
\small
\begin{tabular}{rrrrrrrr}
\toprule
$g$ & HF & MBPT2 & MBPT3 & CID & HF + RPA & CCD & FCI\\
\midrule
$0.25$ & $1.75000$ & $1.73177$ & $1.73047$ & $1.73051$ & $1.71635$
  & $1.7304562$ & $1.7304598$\\
$0.50$ & $1.50000$ & $1.42708$ & $1.41667$ & $1.41760$ & $1.37450$
  & $1.4166377$ & $1.4167743$\\
$1.00$ & $1.00000$ & $0.70833$ & $0.62500$ & $0.64891$ & $0.55459$
  & $0.6304428$ & $0.6355485$\\
$1.50$ & $0.50000$ & $-0.15625$ & $-0.43750$ & $-0.29113$ & $-0.40625$
  & $-0.3841653$ & $-0.3481905$\\
$2.00$ & $0.00000$ & $-1.16667$ & $-1.83333$ & $-1.35209$ & $-1.47622$
  & $-1.6095944$ & $-1.4896522$\\
\midrule
\multicolumn{8}{l}{\emph{errors}}\\
$0.25$ & & $1.3\times10^{-3}$ & $8.9\times10^{-6}$ & $4.7\times10^{-5}$
  & $-1.4\times10^{-2}$ & $-3.6\times10^{-6}$ & \\
$0.50$ & & $1.0\times10^{-2}$ & $-1.1\times10^{-4}$ & $8.2\times10^{-4}$
  & $-4.2\times10^{-2}$ & $-1.4\times10^{-4}$ & \\
$1.00$ & & $7.3\times10^{-2}$ & $-1.1\times10^{-2}$ & $1.3\times10^{-2}$
  & $-8.1\times10^{-2}$ & $-5.1\times10^{-3}$ & \\
$1.50$ & & $1.9\times10^{-1}$ & $-8.9\times10^{-2}$ & $5.7\times10^{-2}$
  & $-5.8\times10^{-2}$ & $-3.6\times10^{-2}$ & \\
$2.00$ & & $3.2\times10^{-1}$ & $-3.4\times10^{-1}$ & $1.4\times10^{-1}$
  & $1.3\times10^{-2}$ & $-1.2\times10^{-1}$ & \\
\bottomrule
\end{tabular}
\caption{The pairing model treated by every method of this book:
Hartree-Fock (Chapter~\ref{chap:hartrefock}), second- and third-order
perturbation theory and doubles configuration interaction
(Chapters~\ref{chap:mbpt} and~\ref{chap:fci}), the random-phase
approximation (Chapter~\ref{chap:mfapplications}), coupled-cluster doubles,
and full configuration interaction as the exact reference.}
\label{tab:cccomparison}
\end{table}

At weak and moderate coupling coupled cluster is the best of the five, by two
to three orders of magnitude at $g=0.25$ and by a factor of two over
third-order perturbation theory at $g=1$.  Unlike the perturbation series it
does not fall apart as the coupling grows: at $g=2$, where MBPT3 is wrong by
$0.34$, CCD is wrong by $0.12$.

It does not win everywhere.  At $g=2$ the random-phase approximation happens
to be closer, which is a reminder that a method summing a different class of
contributions can prevail in the regime it was designed for.  And CCD is not
variational -- the errors change sign across the table -- so a single number
carries no guarantee.

What distinguishes it is that no other entry in the table is simultaneously
size extensive, systematically improvable and accurate across the whole
range.  Configuration interaction is variational but not extensive;
perturbation theory is extensive but divergent; the random-phase
approximation is neither systematically improvable nor stable.  Coupled
cluster is extensive by construction, improvable by raising the rank of
$\hat{T}$, and polynomial in cost.  That combination is why it, and not the
others, became the standard.

%----------------------------------------------------------------------
\section{Beyond CCSD}
\label{sec:beyondccsd}

\index{coupled cluster!triples}\index{CCSD(T)}
The hierarchy continues.  Including $\hat{T}_3$ gives CCSDT, with a third
amplitude equation
\begin{equation}
  0 = \bracket{\Phi_{ijk}^{abc}}{\overline{H}_N}{\Phi_0},
  \label{eq:10-triples}
\end{equation}
whose structure is the same -- a driving term, linear terms, and nonlinear
couplings -- and whose length is formidable.  The cost rises to $n^8$.

Empirically, CCSD recovers around $90\%$ of the correlation energy of a
well-behaved system and CCSDT around $99\%$.  Since the triples are expensive
and their effect modest, the standard compromise is to compute them
perturbatively rather than iteratively.  Approximating
\begin{equation}
  t_{ijk}^{abc} \approx
    \frac{N_{ijk}^{abc}\!\left(t_1,t_2;\langle pq\Vert rs\rangle\right)}
         {D_{ijk}^{abc}},
  \qquad
  D_{ijk}^{abc} = \varepsilon_i+\varepsilon_j+\varepsilon_k
    -\varepsilon_a-\varepsilon_b-\varepsilon_c,
  \label{eq:10-perturbativetriples}
\end{equation}
with a numerator built from the converged singles and doubles, gives
CCSD(T) at a cost of $n^7$ and with no iteration.  This method is accurate
enough, and cheap enough, to have earned the name ``the gold standard of
quantum chemistry''.

Two further directions deserve mention.  Because $\overline{H}_N$ is not
Hermitian, its left and right eigenvectors differ, and computing properties
other than the energy requires both -- this is the price of the exponential
noted in Section~\ref{sec:similaritytransform}.  And diagonalising
$\overline{H}_N$ in the space of excited determinants, rather than merely
demanding that its $1p$-$1h$ and $2p$-$2h$ blocks vanish, gives excitation
energies: this is equation-of-motion coupled cluster, and it stands to the
present chapter as the random-phase approximation of
Chapter~\ref{chap:mfapplications} stands to Hartree-Fock.

%----------------------------------------------------------------------
\section{Summary}
\label{sec:ch10summary}

Coupled-cluster theory replaces the linear ansatz
$(1+\hat{C})\bket{\Phi_0}$ by the exponential $e^{\hat{T}}\bket{\Phi_0}$.
The two are equivalent when complete, and they differ decisively when
truncated: the exponential builds high-rank excitations as products of
low-rank ones, so truncating $\hat{T}$ at doubles retains the disconnected
quadruples that truncating $\hat{C}$ discards.  That single reorganisation
makes the method size extensive at every truncation level, which
configuration interaction is not.

Requiring $e^{-\hat{T}}\hat{H}e^{\hat{T}}$ to have the reference as an
eigenstate turns the Schr\"odinger equation into the projected equations
(\ref{eq:10-energyeq})--(\ref{eq:10-doubleseq}).  Because the cluster
operators commute among themselves, every $\hat{T}$ in the
Baker-Campbell-Hausdorff series must connect to $\hat{H}_N$, which has only
four contractible operators; the expansion therefore terminates exactly at
the fourfold commutator.  The non-Hermiticity that this exploits is the same
non-Hermiticity that complicates properties -- one cannot have the one
without the other.

The resulting equations, (\ref{eq:10-ccd}) for doubles and
(\ref{eq:10-ccsdsingles})--(\ref{eq:10-ccsddoubles}) for singles and doubles,
are nonlinear and are solved by iteration.  Their first iteration is MP2;
their order-by-order expansion is M\o{}ller-Plesset perturbation theory,
verified here through fourth order; and solving them to convergence sums that
expansion within the retained excitation classes.  For the pairing model CCD
is accurate to four parts in a million at weak coupling and still the best
method available at strong coupling, while remaining exactly size extensive.

The program \texttt{coupledcluster.py} performs every calculation of this
chapter.

%----------------------------------------------------------------------
\section{Exercises}
\label{sec:ch10exercises}
%=============================================================================

\subsection*{Exercise 1: the exponential ansatz}
\label{ex:10-ansatz}
\addcontentsline{toc}{subsection}{Exercise 1: the exponential ansatz}

\begin{enumerate}
  \item[a)] Prove Eq.~(\ref{eq:10-Tcommute}), that the cluster operators
  commute, and identify exactly where the particle-hole structure is used.
  \item[b)] Derive the relations~(\ref{eq:10-CtoT}) between the
  configuration-interaction and cluster amplitudes, and extend them to
  $\hat{C}_4$.
  \item[c)] Show that for two non-interacting subsystems the exponential
  ansatz factorises, Eq.~(\ref{eq:10-factorises}), and that the linear ansatz
  truncated at doubles does not.
  \item[d)] Why does the exponential series terminate when acting on an
  $N$-particle determinant?
\end{enumerate}

\paragraph{Answer to a).}
Take two terms of $\hat{T}$, each a string of particle creation and hole
annihilation operators.  Since particle and hole indices are disjoint, every
annihilation operator in one string anticommutes with every creation operator
in the other, with no delta function generated.  Moving one string past the
other therefore produces a sign $(-1)^{n_1n_2}$ from the interchange, and the
same sign appears in the reverse ordering, so the two products are equal.
This is Eq.~(\ref{eq:6-phcommute}) generalised, and it is exactly the
disjointness of the index sets that makes it work.

\paragraph{Answer to b).}
Expand $e^{\hat{T}}$ and collect terms by excitation rank.  Rank one gives
$\hat{C}_1=\hat{T}_1$; rank two, $\hat{C}_2=\hat{T}_2+\tfrac12\hat{T}_1^2$;
rank three, $\hat{C}_3=\hat{T}_3+\hat{T}_1\hat{T}_2+\tfrac16\hat{T}_1^3$; and
rank four,
\begin{equation}
  \hat{C}_4 = \hat{T}_4 + \hat{T}_1\hat{T}_3
    + \tfrac12\hat{T}_2^2 + \tfrac12\hat{T}_1^2\hat{T}_2
    + \tfrac{1}{24}\hat{T}_1^4 .
  \label{eq:10-exb1}
\end{equation}
The term $\tfrac12\hat{T}_2^2$ is the one that matters: it is present in CCD,
where $\hat{T}_3=\hat{T}_4=\hat{T}_1=0$, and absent from CID.

\paragraph{Answer to c).}
With no interaction between the subsystems, no amplitude can carry indices
from both, so $\hat{T}=\hat{T}_A+\hat{T}_B$ with $[\hat{T}_A,\hat{T}_B]=0$.
Hence $e^{\hat{T}}=e^{\hat{T}_A}e^{\hat{T}_B}$ and, since the operators act
on disjoint orbitals, the state factorises.  The energy of a product state
under $\hat{H}_A+\hat{H}_B$ is the sum of the two energies.  For the linear
ansatz, $(1+\hat{C}_2^A+\hat{C}_2^B)$ is \emph{not} the product
$(1+\hat{C}_2^A)(1+\hat{C}_2^B)$: the cross term
$\hat{C}_2^A\hat{C}_2^B$ is a quadruple excitation, which CID has thrown
away.  This is Section~\ref{sec:sizeconsistency} in one line.

\paragraph{Answer to d).}
Each application of $\hat{T}$ annihilates at least one hole, and there are
only $N$ holes.  Hence $\hat{T}^{n}\bket{\Phi_0}=0$ for $n>N$, and more
sharply for a truncated $\hat{T}$: with $\hat{T}=\hat{T}_2$ the series stops
at $\hat{T}_2^{N/2}$.

\subsection*{Exercise 2: the similarity transformation}
\label{ex:10-similarity}
\addcontentsline{toc}{subsection}{Exercise 2: the similarity transformation}

\begin{enumerate}
  \item[a)] Show that $\overline{H}=e^{-\hat{T}}\hat{H}e^{\hat{T}}$ has the
  same eigenvalues as $\hat{H}$ for any $\hat{T}$.
  \item[b)] What condition on $\hat{T}$ would make $\overline{H}$ Hermitian?
  Why does coupled cluster not impose it?
  \item[c)] Prove that the Baker-Campbell-Hausdorff series terminates at the
  fourfold commutator, and say precisely which property of $\hat{H}_N$ is
  used.
  \item[d)] If a similarity transformation preserves all eigenvalues, why is
  the CCD energy not exact?
\end{enumerate}

\paragraph{Answer to a).}
If $\hat{H}\bket{E}=E\bket{E}$ then
\begin{equation}
  \overline{H}\left(e^{-\hat{T}}\bket{E}\right)
   = e^{-\hat{T}}\hat{H}e^{\hat{T}}e^{-\hat{T}}\bket{E}
   = e^{-\hat{T}}\hat{H}\bket{E}
   = E\,e^{-\hat{T}}\bket{E},
  \label{eq:10-exa2}
\end{equation}
so $e^{-\hat{T}}\bket{E}$ is an eigenstate of $\overline{H}$ with the same
eigenvalue.  The map $\bket{E}\mapsto e^{-\hat{T}}\bket{E}$ is invertible, so
the spectra coincide.  Note that eigen\emph{values} are preserved but
eigen\emph{vectors} are not, and neither are norms: $\overline{H}$ is not
Hermitian and its left and right eigenvectors differ, satisfying
$\braket{L_\alpha}{R_\beta}=\delta_{\alpha\beta}$.

\paragraph{Answer to b).}
$\overline{H}$ is Hermitian if $e^{\hat{T}}$ is unitary, that is if
$\hat{T}^{\dagger}=-\hat{T}$.  Coupled cluster does not impose it because an
anti-Hermitian $\hat{T}$ contains de-excitation operators
$\hat{a}^{\dagger}_i\hat{a}_a$ as well as excitations, those do not commute
with the excitations, and the Baker-Campbell-Hausdorff series then never
terminates.  One would gain Hermiticity and lose the finite expansion.
Unitary coupled cluster makes precisely that trade, and is used where the
exponential can be realised directly rather than expanded -- on a quantum
computer, for instance.

\paragraph{Answer to c).}
By Eq.~(\ref{eq:10-Tcommute}) no commutator can be made non-zero by one
$\hat{T}$ acting on another, so in
$[[\cdots[\hat{H}_N,\hat{T}],\ldots],\hat{T}]$ every $\hat{T}$ must be
contracted with $\hat{H}_N$ itself.  Each such contraction consumes one of
$\hat{H}_N$'s operators that can be contracted to the right -- and its
two-body part has exactly four, $\hat{a}^{\dagger}_p\hat{a}^{\dagger}_q$ and
$\hat{a}_s\hat{a}_r$ in normal-ordered form.  A fifth $\hat{T}$ has nothing
left to attach to, so the term vanishes.  The property used is that
$\hat{H}_N$ is \emph{at most two-body}; a three-body force would extend the
series to the sixfold commutator.

\paragraph{Answer to d).}
Because we do not solve the transformed problem exactly either.  With
$\hat{T}=\hat{T}_2$ the full $\overline{H}_N$ contains up to six-body terms,
and it is that operator whose spectrum equals the exact one.  What CCD
imposes is only that the $2p$-$2h$ block of $\overline{H}_N$ vanish -- the
reference is decoupled from doubles but not from triples and quadruples.
Coupled cluster at finite truncation is a similarity transformation
\emph{plus} a projection, and the error lives in the projection.

\subsection*{Exercise 3: deriving the CCD equations}
\label{ex:10-ccd}
\addcontentsline{toc}{subsection}{Exercise 3: deriving the CCD equations}

\begin{enumerate}
  \item[a)] Show that the expansion~(\ref{eq:10-ccdexpansion}) stops at
  $\hat{T}_2^2$ when projected onto a $2p$-$2h$ state.
  \item[b)] Derive the linear terms~(\ref{eq:10-ccd1a})
  and~(\ref{eq:10-ccd1b}) and explain the factors $\tfrac12$ and the
  permutation operators.
  \item[c)] Derive the four quadratic terms~(\ref{eq:10-ccd2}).
  \item[d)] Verify that setting all amplitudes on the right-hand side to zero
  gives the MP2 amplitude.
  \item[e)] Show that the intermediates~(\ref{eq:10-intermediates}) reduce
  the cost of the quadratic terms from $\bigO(n^8)$ to $\bigO(n^6)$.
\end{enumerate}

\paragraph{Answer to a).}
A $2p$-$2h$ projection supplies two hole and two particle external lines.
Each $\hat{T}_2$ must connect to $\hat{H}_N$, consuming two of its
contractible operators, and $\hat{H}_N$ has four.  Two cluster operators
therefore exhaust it, and a third has nothing to attach to.

\paragraph{Answer to b).}
Use the generalised Wick theorem of Section~\ref{sec:wickgeneral} on
$\bracket{\Phi_{ij}^{ab}}{\hat{H}_N\hat{T}_2}{\Phi_0}$.  The one-body part
can contract with either a particle or a hole line of the amplitude, giving
the two terms of Eq.~(\ref{eq:10-ccd1a}); the permutation operators appear
because the external labels $a,b$ and $i,j$ must be antisymmetrised.  For the
two-body part, contracting both of its creation operators with the amplitude's
particle lines gives the particle-particle ladder, both annihilation
operators with the hole lines gives the hole-hole ladder, and one of each
gives the ring.  The factor $\tfrac12$ in each ladder compensates the two
equivalent ways of assigning the summed pair of indices; the ring has
inequivalent lines and instead carries the full $P(ij)P(ab)$.

\paragraph{Answer to c).}
Now both $\hat{T}_2$ connect to $\hat{H}_N$, which forces the two-body part
and uses all four of its operators, so only
$\langle kl\Vert cd\rangle$ contributes.  The four topologies are fixed by
how the four contractions are shared: two with each amplitude in the ``ladder
times ladder'' pattern (first term), one particle and one hole with each
(second term, the ring-ring), or both particle and both hole indices with one
amplitude leaving the other a spectator (third and fourth).  The combinatorial
factors follow by counting equivalent assignments as in b).

\paragraph{Answer to d).}
Setting $t=0$ on the right leaves
$D_{ij}^{ab}t_{ij}^{ab}=\langle ab\Vert ij\rangle$, so
\begin{equation}
  t_{ij}^{ab} = \frac{\langle ab\Vert ij\rangle}
    {\varepsilon_i+\varepsilon_j-\varepsilon_a-\varepsilon_b},
  \label{eq:10-exd3}
\end{equation}
which is Eq.~(\ref{eq:9-firstordercoefficients}).  Substituting into
Eq.~(\ref{eq:10-ccsdenergy}) gives the MP2 energy.  The program checks this
against the closed form
$-\tfrac{g^2}{4}[1/(2+g)+2/(4+g)+1/(6+g)]$ for the pairing model.

\paragraph{Answer to e).}
Written out, a term such as
$\sum_{klcd}\langle kl\Vert cd\rangle t_{ik}^{ac}t_{jl}^{bd}$ has four free
external indices and four summed ones, hence $\bigO(n^8)$.  Building
$\chi_{bkcj}=\tfrac12\sum_{ld}\langle kl\Vert cd\rangle t_{lj}^{db}$ first
costs $\bigO(n^6)$ -- four free, two summed -- and contracting it with the
remaining amplitude costs $\bigO(n^6)$ again.  The saving is the standard one
of never forming an object with more indices than necessary, and it is what
makes the $n^6$ scaling of CCSD achievable.

\subsection*{Exercise 4: CCSD and its perturbative content}
\label{ex:10-ccsd}
\addcontentsline{toc}{subsection}{Exercise 4: CCSD and its perturbative content}

\begin{enumerate}
  \item[a)] Derive the energy expression~(\ref{eq:10-ccsdenergy}) and explain
  why no term with three or more cluster operators appears.
  \item[b)] Show that Eqs.~(\ref{eq:10-ccsdsingles})
  and~(\ref{eq:10-ccsddoubles}) reduce to the CCD equation when
  $\hat{T}_1=0$.
  \item[c)] Show that in a Hartree-Fock basis the singles amplitudes are of
  second order in the interaction while the doubles are of first order.
  \item[d)] Establish the order-by-order recursion~(\ref{eq:10-orderbyorder})
  and show that the first two energies are MP2 and MP3.
  \item[e)] Why does CCD reproduce MP4 for the pairing model but not in
  general?
\end{enumerate}

\paragraph{Answer to a).}
$\bracket{\Phi_0}{\overline{H}_N}{\Phi_0}$ requires that the excitation
produced by the cluster operators be exactly undone by $\hat{H}_N$, which can
de-excite by at most two ranks.  Hence at most $\hat{T}_2$ or
$\tfrac12\hat{T}_1^2$ can appear, plus $\hat{T}_1$ against the one-body part.
Three cluster operators would produce at least a $3p$-$3h$ excitation, which
$\hat{H}_N$ cannot annihilate.

\paragraph{Answer to b).}
With $\hat{T}_1=0$ we have $\tau=\tilde\tau=t_2$; the last two lines of
Eq.~(\ref{eq:10-ccsddoubles}) vanish, as does the second half of the ring
term; and the intermediates collapse to
$F_{ae}\to f'_{ae}-\tfrac12 t_2\!\cdot\!v$,
$W_{mnij}\to\langle mn\Vert ij\rangle+\tfrac14 t_2\!\cdot\!v$ and so on,
which are exactly the combinations of
Eq.~(\ref{eq:10-intermediates}).  What remains is Eq.~(\ref{eq:10-ccd}).

\paragraph{Answer to c).}
In a Hartree-Fock basis $f_{ia}=0$, so the driving term of
Eq.~(\ref{eq:10-ccsdsingles}) vanishes.  The leading surviving term is
$-\tfrac12\sum_{mef}t_{im}^{ef}\langle ma\Vert ef\rangle$, first order in
$t_2$; since $t_2$ is itself first order in the interaction, $t_1$ is second
order.  The doubles, by contrast, are driven directly by
$\langle ij\Vert ab\rangle$ and are first order.  The program shows this
concretely: the singles amplitudes of the pairing plus particle-hole model
drop by an order of magnitude when the reference is rotated to the
Hartree-Fock one.

\paragraph{Answer to d).}
The CCD equation is quadratic, so $D t=C+L[t]+Q[t,t]$ exactly.  Substituting
$t=\sum_k t^{(k)}$ and noting that $C$, $L$ and $Q$ each carry one power of
the interaction, the terms of order $k$ are
$L[t^{(k-1)}]$ and $Q[t^{(i)},t^{(j)}]$ with $i+j=k-1$, which is
Eq.~(\ref{eq:10-orderbyorder}).  The first amplitude is the MP2 one by d) of
the previous exercise, and its energy is MP2.  The second satisfies
$Dt^{(2)}=L[t^{(1)}]$, whose energy contribution is the sum of the
particle-particle ladder, hole-hole ladder and ring contractions of two MP2
amplitudes -- which is precisely the third-order energy.

\paragraph{Answer to e).}
At fourth order the exact energy receives contributions from singles,
doubles, triples and quadruples.  CCD has the doubles, and it has the
quadruples in their disconnected form $\tfrac12\hat{T}_2^2$, which is the
whole of the fourth-order quadruple contribution.  It lacks the singles and
the connected triples.  For the pairing model those are absent anyway -- the
interaction cannot produce a $1p$-$1h$ or a $3p$-$3h$ excitation -- so CCD is
complete at fourth order and the program finds exact agreement.  For a
general interaction the connected triples are present and CCD misses them, so
agreement stops at third order.  This is exactly why CCSD(T) exists.

\subsection*{Exercise 5: numerical explorations}
\label{ex:10-numerical}
\addcontentsline{toc}{subsection}{Exercise 5: numerical explorations}

\begin{enumerate}
  \item[a)] Using \texttt{coupledcluster.py}, reproduce
  Table~\ref{tab:ccdpairing} and extend it to six and eight levels.  Does the
  CCD error grow or shrink with the size of the space?
  \item[b)] Verify the size-extensivity result for four and five subsystems,
  and compare with doubles configuration interaction over the same range.
  \item[c)] Add $\hat{T}_4$ to the pairing calculation, or equivalently solve
  the seniority-zero problem exactly, and confirm that the remaining CCD
  error is entirely due to connected quadruples.
  \item[d)] Study the convergence of the CCD iteration as a function of $g$.
  At what coupling does it stop converging, and how does that compare with
  the radius of convergence of the perturbation series in
  Section~\ref{sec:convergence}?
  \item[e)] Implement the CCSD(T) correction of
  Eq.~(\ref{eq:10-perturbativetriples}) for the pairing plus particle-hole
  model and see how much of the remaining error it removes.
\end{enumerate}

\paragraph{Answer.}
Some pointers.  In a) the relative error shrinks as the space grows, because
the neglected connected quadruples become a smaller fraction of a larger
correlation energy -- the opposite of the trend for truncated configuration
interaction, whose size-consistency error grows.  In b) the coupled-cluster
error stays at machine precision for any number of subsystems, while the CID
error grows roughly linearly.  In c) the pairing model with four levels has
no $6p$-$6h$ states, so CCD plus connected quadruples is exact, and the
difference from CCD is precisely the number in the error column of
Table~\ref{tab:ccdpairing}.  In d) the CCD iteration continues to converge
well past the point where the perturbation series diverges, which is the
practical meaning of resummation; it eventually fails too, and the failure is
associated with the same physics that made the mean field unstable in
Chapter~\ref{chap:hartrefock}.  In e) the perturbative triples remove a large
part of the residual error at weak coupling and much less of it at strong
coupling, where the expansion underlying them is itself unreliable.
